\section{Stability and spectral diagnostics for the barycentric field}
\label{sec:asymptotics}

Two questions organize this appendix: what makes the constructed field a stable
peer-averaging system, and how close is the fitted model to the
strong-interaction boundary?
We first state the maintained conditions and derive their implications for peer
averaging, stationary weighting, and the spatial multiplier.
We then report the fitted residual, Dobrushin, and finite-$\rho$ diagnostics
before characterizing the rank-one boundary.
These results certify the constructed field, although their matrix implications
apply to any row-stochastic interaction matrix that satisfies the stated
conditions.

\subsection{Maintained operator conditions}

The first question is which matrix properties are needed for this
interpretation.
The main argument requires a nonnegative row-stochastic interaction matrix
$W$, a strictly positive stationary probability vector $\pi$, and a coefficient
$\rho$ that satisfies the stability condition in \Cref{hyp:spectral-gate}.
Row stochasticity defines the directed peer average, the stationary vector
supplies the weighting in \Cref{thm:spatial-attenuation}, and the stability
condition makes the spatial multiplier well defined.
Let $P=\1\pi^\top$.
All matrix norms here are maximum absolute row-sum norms.

In the application, $W=W^\flat$ is the barycentric interaction field obtained
from target-anchored reconstruction.
The spectral results operate on that fitted matrix.

The boundary analysis also requires a rate at which repeated peer averaging
approaches $P$.
The following certificate states that condition.

\begin{definition}[Geometric Perron certificate]
  \label{def:perron-limit}
  A row-stochastic matrix $W$ has a geometric Perron certificate
  $(\pi,C,q)$ if $0\leq q<1$ and
  \[
    \lVert W^k-P\rVert \leq Cq^k
    \qquad\text{for every }k\geq0.
  \]
\end{definition}

\subsection{Implications for peer averaging and the multiplier}

The next question is whether the maintained conditions produce coherent
long-run peer weights.
Row stochasticity makes $Wx$ a peer average, while the stability condition
makes $(I-\rho W)^{-1}$, and hence the equilibrium mapping from stand-alone
exposures, well defined.
The geometric certificate adds the limiting result needed for the boundary
analysis.
\begin{theorem}[Geometric certificate implies the Perron limit]
  \label{thm:perron-from-geometric}
  If $W$ has a geometric Perron certificate $(\pi,C,q)$, then
  $W^k\to P$ as $k\to\infty$.
\end{theorem}
Because $0\leq q<1$, $Cq^k\to0$, so the defining inequality yields
$W^k\to P$.

The Perron limit has two consequences for the economic aggregation in the main
text.
First, repeated peer averaging replaces any initial profile by its
$\pi$-weighted cross-sectional mean.
\begin{theorem}[Stationary interaction centrality]
  \label{thm:replication-centrality}
  Under the Perron limit, $W^k x\to(\pi^\top x)\1$ for every vector $x$.
\end{theorem}
Multiplying $W^k\to P=\1\pi^\top$ by a fixed vector $x$ gives
$W^k x\to Px=(\pi^\top x)\1$.

Second, the weights that define this long-run mean are invariant under one
application of the operator.
\begin{theorem}[Stationarity]\label{thm:stationarity}
  Under the Perron limit, $\pi^\top W=\pi^\top$.
\end{theorem}
Since $W^{k+1}=W^k W$, the Perron limit gives $W^{k+1}\to PW$, while the
same limit with $k+1$ gives $W^{k+1}\to P$. Uniqueness of limits implies
$PW=P$, and $P=\1\pi^\top$ then gives $\pi^\top W=\pi^\top$.

Together, the two results make $\pi_i$ firm $i$'s long-run influence and justify
the stationary centring in \Cref{thm:spatial-attenuation}.

\subsection{Fitted operator and finite-\texorpdfstring{$\rho$}{rho} diagnostics}

The empirical diagnostic question is whether the fitted field and coefficient
satisfy these conditions and whether the boundary closely approximates the
empirical model.
The fitted $\ppvalue{n-tickers-priced}\times\ppvalue{n-tickers-priced}$
operator is row stochastic to a maximum residual of
\ppnum[8]{network-row-sum-residual}.
Because its entries are nonnegative and its rows sum to one, its maximum
absolute row-sum norm is one.
The pooled estimate
$\hat\rho=\ppnum[3]{post-pooled-2023-2026-w-flat-rho}$ therefore satisfies the
stability condition and makes the spatial multiplier well defined.
Power iteration also yields a strictly positive stationary distribution with
minimum mass \ppnum[4]{network-stationary-minimum} and stationarity residual
\ppnume[2]{network-stationary-residual}.
These fitted checks support, respectively, the peer average, the equilibrium
mapping from stand-alone exposures, and the stationary weighting used in the
main text.

The Dobrushin coefficient
$\delta(W)=\tfrac12\max_{i,j}\sum_m|W_{im}-W_{jm}|$ measures the largest
difference between any two firms' peer-weight profiles. A value below one means
that all row pairs retain some overlap and supplies the auditable envelope
$C=2$ and $q=\delta(W)$: every row of $W^k$ lies within
$2\delta(W)^k$ of the stationary distribution in $\ell_1$ distance.
For the fitted operator, the second power is strictly positive with minimum
entry \ppnum[6]{network-w2-minimum-weight}, and the Dobrushin coefficient is
\ppnum[3]{network-dobrushin}.
Thus the fitted operator has a geometric Perron certificate: its iterated peer
profiles converge even though the operator is directed.

The Perron limit describes the boundary $\rho\uparrow1$, whereas the empirical
application uses a finite working-model coefficient.
An exact decomposition separates the rank-one part of the multiplier from the
remaining cross-sectional variation at any admissible $\rho$.
Write $N=W-P$. Row stochasticity, stationarity, and unit mass imply
$WP=PW=P^2=P$. Whenever the displayed inverses exist and $\rho\neq1$, the
Leontief multiplier splits exactly as
\begin{equation}
  (I-\rho W)^{-1}
  =(1-\rho)^{-1}P+(I-\rho N)^{-1}(I-P).
  \label{eq:resolvent-split}
\end{equation}
The error made by replacing the normalized multiplier with its limiting
projection is therefore explicit rather than asymptotic:
\begin{equation}
  (1-\rho)(I-\rho W)^{-1}-P
  =(1-\rho)(I-\rho N)^{-1}(I-P).
  \label{eq:finite-rho-remainder}
\end{equation}

\begin{proposition}[Finite-$\rho$ resolvent bound]
  \label{prop:finite-rho-bound}
  If $\lVert\rho(W-P)\rVert<1$, then
  \[
    \left\lVert(1-\rho)(I-\rho W)^{-1}-P\right\rVert
    \leq
    |1-\rho|\{1-\lVert\rho(W-P)\rVert\}^{-1}\lVert I-P\rVert.
  \]
\end{proposition}
The proposition provides an auditable upper bound on the distance between the
finite-$\rho$ normalized multiplier and its limiting projection.
Taking norms in \eqref{eq:finite-rho-remainder} and applying the Neumann bound
\[
  \lVert(I-\rho N)^{-1}\rVert
  \leq\{1-\lVert\rho N\rVert\}^{-1},
\]
gives the displayed bound.

At the pooled fitted coefficient, the exact normalized-resolvent error is
\ppnum[3]{network-w-flat-resolvent-error}, against a Dobrushin bound of
\ppnum[3]{network-w-flat-resolvent-bound}.
The Neumann envelope does not apply on this panel: it requires
$\lVert\rho N\rVert<1$, and the fitted coefficient gives
$\lVert\rho N\rVert=\ppnum[3]{network-w-flat-complement-norm}$, outside the
region of convergence.
Every fitted variant lies outside it on the
\ppnum[0]{n-tickers-priced}-firm panel, so the Dobrushin bound is the only
envelope reported here.
The exact error remains materially above zero, so the rank-one limit is not a
close approximation to the fitted normalized multiplier in this norm.
The Dobrushin bound is a valid but conservative envelope.
Thus the fitted coefficient lies inside the stable region, but the boundary
approximation is not the empirical working model.

\subsection{The rank-one boundary}

The boundary question is what remains of cross-sectional variation as peer
adjustment approaches its strong-interaction limit.
The boundary result follows in two steps.
First, convergence of repeated peer averaging makes the normalized spatial
multiplier converge to the stationary projection.
\begin{theorem}[Resolvent limit from the power limit]
  \label{thm:resolvent-from-perron}
  Let $W$ be row stochastic and suppose $W^k\to P$. Then
  $(1-\rho)(I-\rho W)^{-1}\to P$ as $\rho\uparrow1$.
\end{theorem}
This result converts the power limit into a statement about spatial feedback.
The normalized Neumann series is
\[
  S_\rho=(1-\rho)(I-\rho W)^{-1}
  =(1-\rho)\sum_{k\geq0}\rho^k W^k.
\]
For any $\varepsilon>0$, choose $K$ so that
$\lVert W^k-P\rVert<\varepsilon$ for $k\geq K$. Since
$P=(1-\rho)\sum_{k\geq0}\rho^k P$, split the difference into its finite head
and tail:
\[
  \lVert S_\rho-P\rVert
  \leq(1-\rho)\sum_{k<K}\rho^k\lVert W^k-P\rVert
  +(1-\rho)\sum_{k\geq K}\rho^k\lVert W^k-P\rVert.
\]
The finite head vanishes as $\rho\uparrow1$, while the tail is at most
$\varepsilon$; the uniform bound
$\lVert W^k-P\rVert\leq2$ justifies the split. Hence
$S_\rho\to P$.

Second, substituting that multiplier limit into the spatial covariance removes
all cross-sectional directions except the common one.
\begin{theorem}[Rank-one collapse of the rescaled spatial covariance]
  \label{thm:rankone-collapse}
  Let $\Sigma_{\mathrm{SAR}}=(I-\rho W)^{-1}V(I-\rho W)^{-\top}$. Under the
  Perron limit,
  \[
    (1-\rho)^2\Sigma_{\mathrm{SAR}}
    \longrightarrow (\pi^\top V\pi)\1\1^\top
    \quad\text{as }\rho\uparrow1.
  \]
\end{theorem}
The proof follows by continuity of matrix multiplication.
Set $S_\rho=(1-\rho)(I-\rho W)^{-1}\to P$. Continuity of matrix
multiplication gives
$S_\rho V S_\rho^\top\to PVP^\top$. Since $P=\1\pi^\top$,
\[
  PVP^\top=(\pi^\top V\pi)\1\1^\top.
\]

\Cref{thm:rankone-collapse} is the $\rho\uparrow1$ endpoint of the attenuation
result: at the boundary, spatial feedback removes cross-sectional dispersion
entirely and the covariance becomes rank one.
The finite-$\rho$ remainder shows why this endpoint is useful as a theoretical
boundary but not as an approximation to the fitted working model.
